% These sections report archived, verified results; no new inference is computed.
\subsection{需求感知搜索的直接增量与负载分解}

事后补充S1以DA减FHS5直接评价搜索增量，沿用原优化初态及留出请求。每相位、每规模60个父图，7条轨迹先在父图内平均，3模型等权；20\,000次模型分层配对重抽样对两相位、4规模、2终点的16项比较作Bonferroni分配，单区间名义水平为99.6875\%。全部校正区间均包含或接触0，未支持该比较家族下的严格增益，亦不证明等效。30节点时间差在正式、确认相位为0.001858和0.000906；240节点两终点差均为0（补充表S1、图S1）。

补充S2按同分布4条、偏移3条轨迹描述分解，总体差满足$D=(4D_{\mathrm{same}}+3D_{\mathrm{shift}})/7$。正式相位120节点时间差分别为$-0.000214$与0.000089，合并为$-0.000084$；风险差分别为0.004167与$-0.005556$，合并为0。相反方向可抵消，故合并差不代表每种负载均改善；此分解未作新的范围交互检验。

% SERVICE-COST
补充S3进一步按成功请求归一协调暴露量。令S为完整窗口内成功请求数，C为相应累计路径暴露量，采用
\begin{equation}
R_C=\frac{\mathbb E[C]}{\mathbb E[S]}.
\label{eq:s3-ratio}
\end{equation}
期望按父图内轨迹平均及模型等权定义；二元上下参照先在每条轨迹内平均。式\ref{eq:s3-ratio}是均值之比，不是$\mathbb E[C/S]$或单位支付金额成本。首次无路径后继续处理的请求仍计入S、C，因此$S/H$不同于首次失败风险。确认相位240节点DA及其二元参照的成功比例分别为0.999918、0.962056；每成功请求经过超边数为2.478775、4.379770，信令槽位为12.025870、8.759540，二次暴露量为59.109145、17.519080。短路径与较高协调负担可同时存在（补充图S2、S3）；这些量不是实际时延或费用。比值差由同一父图重抽样内重新求比后相减，区间为未校正点态95\%。

% WINDOW
\subsection{观测窗口敏感性}

补充S4复用原事件记录，将$H=12n$截短为$h=6n$。对$T=\min(\tau_{\mathrm{np}},H)$、$\delta=\mathbf1\{\tau_{\mathrm{np}}\leq H\}$，先在各臂计算$T_h=\min(T,h)$、$\delta_h=\delta\mathbf1\{T\leq h\}$及$Y_h=T_h/h$，再配对；在第h次发生事件与h内无事件的$Y_h$可相同，但$\delta_h$不同。两相位、4规模、4构造和两窗口的64个条件均保持时间差正、风险差负；32组短窗口时间增量均较小，风险差均向0收缩，这只是方向计数，不是独立重复。DA减FHS5的短窗口时间差在8个条件中7个为0，不能据此证明等效。窗口变化用共享索引直接计算，240个区间均为未校正点态95\%；该分析不估计未删失寿命、纯规模效应或短窗口成本（补充图S4、S5）。
